% RESS abstract (<= 200 words; checked by tools/check_front_ress.py). Numbers via numbers.tex.
A perception system on a new fixed camera is often accepted by a zero-failure reliability demonstration: a fixed number
of real events, the zero-failure sample size (\nNzero{} events for 95\% confidence that at most 5\% of events are
missed), must pass without a miss. How many can simulated evidence replace? We prove that,
without an assumption linking the simulator to reality, independent simulated data cannot shorten the demonstration;
that accepting a larger error probability buys an exactly computable number of events; and that a synthetic twin paired
with each real event bounds the real miss rate by the twin miss rate plus a discordance rate, which pays off only
through a transfer bound priced in calibration scenes. On \nMainScenes{} fixed-camera scenes, real misses are rare at the pre-specified operating point (\nPmainOPshort), so
it is untested there. If calibration carries over, a twin run about \nTwinRunsOP{} times per camera would save at most \nSavedOP{} events per camera at a 20\% target; the twin as built saves none. The transfer certificate stands at scene level but is withdrawn at site level (\nSiteSm{} sites). Two named failure modes: night and
camera-jitter cameras pass through (\nOODSNightfa{} of \nOODSNightcams{} and \nOODSJitfa{} of \nOODSJitcams{} falsely accepted).
